\documentclass[10pt,a4paper]{article}
\usepackage[margin=2.2cm]{geometry}
\usepackage{amsmath,amssymb,bm}
\usepackage{booktabs}
\usepackage{mathtools}
\usepackage{caption}
\usepackage{amsmath}
\usepackage{graphicx}
\usepackage{microtype}
\usepackage[colorlinks=true,linkcolor=nvy,citecolor=nvy,urlcolor=nvy]{hyperref}
\usepackage{xcolor}
\definecolor{nvy}{HTML}{0F2440}
\definecolor{red2}{HTML}{B03A2E}
\usepackage{titlesec}
\titleformat{\section}{\normalfont\large\bfseries\color{nvy}}{\thesection}{0.6em}{}
\titleformat{\subsection}{\normalfont\bfseries}{\thesubsection}{0.6em}{}
\newcommand{\ket}[1]{\lvert #1 \rangle}
\newcommand{\bra}[1]{\langle #1 \rvert}
\newcommand{\ip}[2]{\langle #1 | #2 \rangle}
\newcommand{\ssig}{\bm{\sigma}}
\newcommand{\note}[1]{\textcolor{red2}{#1}}
\setlength{\parskip}{0.35em}

\title{\vspace{-1.6cm}\textbf{Variational ground states of three spin-chain
Hamiltonians}\\[2pt]
\large XXX Heisenberg, Haldane--Shastry, and Sachdev--Ye--Kitaev}
\author{Boris-Alexandre Baudel}
\date{}

\begin{document}
\maketitle
\vspace{-1.0cm}

\section*{What is in here, and what is not}

Three models, three different answers, and I think the differences are the
interesting part of the exercise.

\textbf{Model 1 (XXX).} Solved. A Hamiltonian Variational Ansatz that is
exactly implementable (no Trotter error) and that is protected by $SU(2)$
symmetry reaches machine precision for $L \le 6$ and $10^{-6}$ relative error
at $L=8$. At $L=12,16$ a random initialisation stalls at a
few percent; growing the circuit one layer at a time from the previous optimum
recovers a factor of $10$ at $L=12$ and a factor of $105$ at $L=16$
($1.6\times10^{-2} \to 1.5\times10^{-4}$). Both strategies are reported,
because the difference between them is the actual lesson.

\textbf{Model 2 (Haldane--Shastry).} Solved analytically, and this is the
honest answer rather than a variational one. Three lines of algebra reduce the
Hamiltonian as written to the textbook Haldane--Shastry chain; its ground state
is then the Gutzwiller--Jastrow wavefunction, which I construct explicitly and
verify to be an exact eigenvector with unit overlap with the numerical ground
state at every accessible $L$. A VQE here would be a variational
approximation to something already known in closed form.

\textbf{Model 3 (SYK).} Not solved --- and the report explains \emph{why}
with measured numbers rather than adjectives. Five obstructions are quantified.
Two are engineering problems with known partial cures; one --- the volume-law
entanglement of the ground state --- is a statement about circuit expressivity
that no optimiser tuning can fix; and one turned out to be embarrassingly
simple and is usually not on the list, namely that the ground state is
intrinsically \emph{complex} while the textbook hardware-efficient ansatz has
real amplitudes only. ADAPT-VQE is implemented rather than merely
recommended, and the result is negative: with a Hamiltonian-derived pool it
beats the real-amplitude brickwork by a mile but does \emph{not} beat a
complex-capable one at these sizes. Reporting that seemed more useful than
tuning the budget until the ordering flipped.

\textbf{Two things that are usually asserted and are measured here instead.}
Sections~\ref{sec:shots} and~\ref{sec:warm}: the sampling cost is
measured with an unbiased grouped estimator rather than bounded by a
formula ($\sim\!10^8$ shots for one SYK energy at $N=12$ to $10^{-3}$, and the
ordering across models is set by $\|c\|_1$, not by the number of terms); and
the Gaussian $q=2$ warm start is implemented rather than proposed --- it is
what turns the negative ADAPT result positive.

Everything numerical is cross-validated: the fast in-house state-vector engine
is checked against Qiskit, the analytic gradient against finite differences,
the closed-form energies against exact diagonalisation, and the reduced
$S^z$ sector against the full Hilbert space. All numbers in the tables are
generated from the run outputs, not typed in.
Appendix~\ref{app:math} derives every identity used anywhere in the report and
\emph{evaluates} it: $47$ checks, $47$ passing, each with its residual listed.
One of them is a trap worth naming: the parameter-shift rule for the exchange
gate does exist, but only in rescaled form, and it stops being usable for a
whole layer --- the appendix gives both statements and the frequency count
behind them.

\section{Conventions}

\paragraph{Pauli versus spin operators.} The assignment writes the
Hamiltonians with Pauli matrices. Since $\ssig = 2\bm{S}$,
\begin{equation}
  \ssig_i \cdot \ssig_j = 4\, \bm{S}_i \cdot \bm{S}_j ,
\end{equation}
and every comparison with a textbook result has to carry that factor of four.
It is the single most common way to be wrong by a factor of $4$ in this
problem, so it is stated once and used consistently.

\paragraph{Qubit ordering.} Internally site $0$ is the most significant tensor
factor, $\ket{s_0 s_1 \dots s_{L-1}}$. Qiskit is little-endian. All Pauli
labels handed to Qiskit are therefore reversed; the cross-check in
Section~\ref{sec:soft} verifies that the two conventions give identical
expectation values.

\paragraph{Sectors.} All three Hamiltonians commute with $S^z_{\rm tot}$, and
Models 1 and 2 also commute with $\bm{S}^2_{\rm tot}$. Exact diagonalisation is
performed in the $S^z = 0$ sector ($\binom{16}{8} = 12870$ states instead of
$65536$ at $L=16$) and cross-checked against the full space wherever that is
affordable.

\section{Model 1: the XXX Heisenberg chain}

\subsection{Benchmark}

For $H_{\rm XXX} = J\sum_{i=1}^{L}\ssig_i\cdot\ssig_{i+1}$ with periodic
boundary conditions, the Bethe-ansatz result for the isotropic
antiferromagnet, $E/L \to J(\tfrac14 - \ln 2)$ in the spin normalisation,
becomes
\begin{equation}
  \frac{E_0}{L} \;\longrightarrow\; 4J\left(\tfrac14 - \ln 2\right)
  = J\,(1 - 4\ln 2) = -1.7725887\,J .
\end{equation}
This is the target for the finite-size extrapolation.

\subsection{Ansatz: why this one}

The PBC chain with even $L$ splits into two sets of \emph{disjoint} bonds,
\begin{equation}
  B_{\rm odd} = \{(0,1),(2,3),\dots\},\qquad
  B_{\rm even} = \{(1,2),(3,4),\dots,(L{-}1,0)\},
\end{equation}
so that $H = H_{\rm odd} + H_{\rm even}$ with each $H_\bullet$ a sum of
commuting two-site terms. The Hamiltonian Variational Ansatz
\begin{equation}
  \ket{\psi(\bm\theta)} = \prod_{p=1}^{P}
  e^{-i\theta_{p,2} H_{\rm even}}\, e^{-i\theta_{p,1} H_{\rm odd}}
  \ket{\psi_0},
  \qquad
  \ket{\psi_0} = \bigotimes_{(a,b)\in B_{\rm odd}}
  \frac{\ket{01}-\ket{10}}{\sqrt 2},
  \label{eq:hva}
\end{equation}
has three properties that matter.

\emph{(i) It is exact, not Trotterised.} Because the bonds within a layer are
disjoint, $e^{-i\theta H_{\rm odd}}$ factorises \emph{exactly} into two-site
gates. There is no Trotter error to disentangle from the variational error.

\emph{(ii) The two-site gate is elementary.} $\ssig_a\cdot\ssig_b$ has
eigenvalue $+1$ on the triplet and $-3$ on the singlet, i.e.\ it is the
identity on $\{\ket{00},\ket{11}\}$ and $-\mathbb{1}+2X$ on
$\{\ket{01},\ket{10}\}$. Hence
\begin{equation}
  e^{-i\theta\,\ssig_a\cdot\ssig_b} =
  \begin{cases}
    e^{-i\theta} & \text{on } \ket{00},\ket{11},\\[2pt]
    e^{+i\theta}\big(\cos 2\theta\,\mathbb{1} - i \sin 2\theta\, X\big)
      & \text{on } \ket{01},\ket{10}.
  \end{cases}
  \label{eq:gate}
\end{equation}
Exploiting this block structure instead of a generic $4\times4$ contraction is
what makes $L=16$ run in seconds. In circuit form, since $XX$, $YY$ and $ZZ$
commute pairwise on a given pair,
\begin{equation}
  e^{-i\theta(XX+YY+ZZ)} = R_{XX}(2\theta)\,R_{YY}(2\theta)\,R_{ZZ}(2\theta),
\end{equation}
which is the decomposition used for the Qiskit circuits.

\emph{(iii) It cannot leave the right symmetry sector.} Every gate
$e^{-i\theta\ssig\cdot\ssig}$ is an $SU(2)$ scalar and the reference state is a
product of singlets, so the whole ansatz stays inside the total-singlet,
$S^z=0$ sector --- exactly where the antiferromagnetic ground state lives.
The variational search therefore never wastes expressivity on states that are
orthogonal to the target by symmetry. This is the main reason the ansatz is
usable at $16$ qubits with only $2P$ parameters.

\subsection{Gradients}

For a sub-layer $U_k = e^{-i\theta_k H_k}$ with $H_k$ a sum of disjoint terms,
$\partial_{\theta_k} U_k = -i H_k U_k$, hence
\begin{equation}
  \frac{\partial E}{\partial \theta_k}
  = 2\,\mathrm{Re}\,\bra{\lambda}(-i H_k)\ket{\psi_{k+1}},
  \qquad
  \ket{\lambda} = \Big(\textstyle\prod_{j>k} U_j\Big)^{\dagger} H \ket{\psi},
  \label{eq:adjoint}
\end{equation}
with $\ket{\psi_{k+1}}$ the state after sub-layer $k$. Reverse-mode (adjoint)
accumulation gives the full gradient for the cost of about two forward passes,
against $2\times$(number of parameters) for finite differences. Measured
agreement between the two: $\le 3\times10^{-9}$ at $L=16$; measured speed-up at
$L=16$, $P=8$: $\times 8$. A word on parameter-shift rules, because the
natural guess is wrong in both directions. The \emph{naive} $\pi/2$-shift rule
does not apply to \eqref{eq:gate} --- $\ssig\cdot\ssig$ has eigenvalues $+1$
and $-3$, so the single frequency of $E(\theta)$ is $\Omega=4$, not $2$ --- but
a rescaled two-term rule,
$\partial_\theta E = 2[E(\theta+\pi/8)-E(\theta-\pi/8)]$, \emph{is} exact for
one bond. For a whole sub-layer of $R$ disjoint bonds the frequencies are
$4,8,\dots,4R$ and an exact rule needs $2R$ evaluations; at $L=8$ the two-term
rule is then wrong by $2.2$ in absolute terms. Appendix~\ref{app:math}
(\textsf{B.5}, \textsf{B.6}) derives and checks both statements. That count is
why the code differentiates by \eqref{eq:adjoint} instead.

\subsection{Results}

\begin{center}\small
\input{tables/model1_main}
\end{center}

\begin{figure}[h]\centering
\includegraphics[width=\textwidth]{figures/model1_error_vs_layers.pdf}
\caption{Relative error of the HVA against exact diagonalisation.
\emph{Left:} random initialisation --- exact for $L\le 6$, but stalling near a
few percent for $L \ge 12$, with extra layers buying nothing. \emph{Right:}
layerwise warm start, where the $(P{+}1)$-layer circuit is initialised from the
converged $P$-layer solution plus a near-identity layer. The stall is a local
minimum of the optimiser, not a limitation of the ansatz.}
\end{figure}

The single most useful observation of this section is the difference between
the two panels. With random initialisation the error curve flattens at
$\sim 2\times10^{-2}$ for $L=12$ and $16$ and is insensitive to $P$ --- the
signature of a local minimum, not of an inexpressive ansatz. Growing the
circuit one layer at a time, each time starting from the previous optimum with
$\theta_{\rm new}\approx 0$, improves $L=8$ from $1.7\times 10^{-3}$ to
$1.3\times 10^{-6}$ and $L=12$ by an order of magnitude. This is the discrete
analogue of adiabatic state preparation and costs nothing extra.

\begin{figure}[h]\centering
\includegraphics[width=0.46\textwidth]{figures/model1_finite_size.pdf}
\caption{Finite-size scaling of the exact ground-state energy per site against
$1/L^2$, compared with the Bethe-ansatz value $1-4\ln 2$. A two-parameter fit
on five sizes lands within $3\times10^{-3}$; the residual is expected, since
the critical chain carries logarithmic corrections that a pure $1/L^2$ form
does not capture.}
\end{figure}

\noindent The extrapolation gives, with
$(\text{fit}, \text{Bethe}, |\text{deviation}|) = ($\input{tables/model1_extrap}$)$.

\section{Model 2: the Haldane--Shastry chain}

\subsection{The Hamiltonian as written is the textbook one}

With $z_j = e^{2\pi i j/L}$ on the unit circle,
\begin{align}
  |z_j - z_k|^2
  &= \Big| e^{2\pi i j/L} - e^{2\pi i k/L} \Big|^2
   = \Big| e^{\pi i (j+k)/L} \Big|^2
     \Big| e^{\pi i (j-k)/L} - e^{-\pi i (j-k)/L} \Big|^2 \nonumber\\[2pt]
  &= \Big| 2 i \sin\!\big(\pi (j-k)/L\big) \Big|^2
   = 4 \sin^2\!\Big(\frac{\pi (j-k)}{L}\Big).
  \label{eq:sinid}
\end{align}
Substituting \eqref{eq:sinid} and $\ssig_j\cdot\ssig_k = 4\bm S_j\cdot\bm S_k$
into the definition, the two factors of $4$ cancel:
\begin{equation}
  H_{\rm HS}
  = \frac{\pi^2 J}{L^2}\sum_{j<k}
    \frac{\ssig_j\cdot\ssig_k}{|z_j-z_k|^2}
  = J\left(\frac{\pi}{L}\right)^{\!2}
    \sum_{j<k} \frac{\bm S_j \cdot \bm S_k}{\sin^2\!\big(\pi (j-k)/L\big)} .
  \label{eq:hsstd}
\end{equation}
The right-hand side is the Haldane--Shastry chain in its standard
normalisation. The geometric dressing in the problem statement is therefore
cosmetic --- but only once \eqref{eq:sinid} is in hand, and
\eqref{eq:sinid} is checked numerically operator-by-operator (last column but
two of the table below) rather than asserted.

\subsection{Closed-form energy, and the exact ground state}

Given \eqref{eq:hsstd}, the known Haldane--Shastry ground-state energy applies
verbatim,
\begin{equation}
  E_0 = -\frac{\pi^2 J}{24}\left( L + \frac{5}{L} \right),
  \qquad \frac{E_0}{L} \xrightarrow[L\to\infty]{} -\frac{\pi^2}{24}
  = -0.4112335 ,
  \label{eq:hsE}
\end{equation}
and the ground state is the Gutzwiller--Jastrow wavefunction. In the
$S^z=0$ sector, with down spins at sites $n_1 < \dots < n_M$, $M = L/2$,
\begin{equation}
  \psi_{\rm J}(n_1,\dots,n_M) \;\propto\;
  \prod_{i=1}^{M} z_{n_i} \prod_{i<j} \big( z_{n_i} - z_{n_j} \big)^2 .
  \label{eq:jastrow}
\end{equation}
Both \eqref{eq:hsE} and \eqref{eq:jastrow} are \emph{verified, not assumed}:

\begin{center}\small
\input{tables/model2_main}
\end{center}

The closed-form energy agrees with exact diagonalisation to
$\sim10^{-15}$ at every $L$; the state \eqref{eq:jastrow} has eigenvector
residual $\|H\psi - E\psi\| \sim 10^{-15}$ and overlap with the numerical
ground state equal to $1$ to twelve digits. The momentum factor
$\prod_i z_{n_i}$ matters: without it the Jastrow form is still an exact
eigenvector, but of an excited state --- its overlap with the ground state is
zero. That is the kind of detail that is invisible analytically and obvious
numerically, which is why both were run.

\begin{figure}[h]\centering
\includegraphics[width=\textwidth]{figures/model2_verification.pdf}
\caption{\emph{Left:} the three independent checks --- the geometric identity,
the closed-form energy and the Jastrow eigenvector residual --- all sitting at
machine precision. \emph{Right:} $E_0/L$ against $1/L^2$, following
\eqref{eq:hsE} exactly, with no fitting.}
\end{figure}

\subsection{So what should a VQE do here?}

Nothing, and saying so is the point. Model 2 has a zero-parameter exact
ansatz. Running a variational circuit on it would produce a controlled
approximation to a quantity already known in closed form, and the only honest
uses of \eqref{eq:jastrow} in a variational setting are:

\begin{itemize}\itemsep2pt
\item as a \textbf{warm start} --- prepare $\psi_{\rm J}$ and use it as the
  reference state of a short HVA built from the long-range bonds, so that the
  circuit only has to represent corrections. Applied to a perturbed
  Hamiltonian (e.g.\ $1/r^2 \to 1/r^{2+\epsilon}$, which destroys
  integrability) this is a genuinely useful construction, and it is the reason
  the exact solution is worth having.
\item as a \textbf{benchmark of a generic ansatz}: the nearest-neighbour HVA
  of \eqref{eq:hva} applied to the all-to-all $H_{\rm HS}$ is a legitimate test
  of what happens when the ansatz structure does not match the interaction
  graph --- the opposite regime from Model 1.
\end{itemize}

The practical cost also changes: $H_{\rm HS}$ has $3\binom{L}{2}$ Pauli terms
against $3L$ for the XXX chain ($84$ versus $24$ at $L=8$, from the run log),
so even measuring the energy is a different problem.

\section{Model 3: the SYK model}

\subsection{Jordan--Wigner mapping}

With $N$ Majoranas and $n = N/2$ qubits, take
\begin{equation}
  \chi_{2k} = \Big(\prod_{j<k} Z_j\Big) X_k,
  \qquad
  \chi_{2k+1} = \Big(\prod_{j<k} Z_j\Big) Y_k,
  \qquad k = 0,\dots,n-1 .
  \label{eq:jw}
\end{equation}
The algebra $\{\chi_a,\chi_b\} = 2\delta_{ab}$ follows from two observations:
on the same qubit, $\{X_k,Y_k\}=0$ and $X_k^2=Y_k^2=\mathbb{1}$; on different
qubits $k<l$, the string of $\chi$ with the larger index contains $Z_k$, and
$Z_k$ anticommutes with $X_k$ and $Y_k$, so the two operators anticommute.
This is verified numerically for $N = 4,\dots,12$ by explicit matrix
construction.

Because each $\chi_a$ is a single Pauli string, so is every quartic term:
$\chi_i\chi_j\chi_k\chi_l$ is one Pauli string times a phase in
$\{\pm1,\pm i\}$. Also, reversing a product of four distinct anticommuting
factors costs $(-1)^{\binom{4}{2}} = +1$, so each term is Hermitian and $H$
is Hermitian for real couplings, as required. With $J_{ijkl}$ totally
antisymmetric the $1/4!$ of the problem statement exactly cancels the $4!$
orderings, so
\begin{equation}
  H_{\rm SYK} = \sum_{i<j<k<l} J_{ijkl}\, \chi_i\chi_j\chi_k\chi_l ,
  \qquad
  \overline{J_{ijkl}} = 0,
  \quad \overline{J^2_{ijkl}} = \frac{6J^2}{N^3} .
\end{equation}

\subsection{Four obstructions, measured}

\begin{center}\small
\input{tables/model3_counting}
\end{center}

\paragraph{(a) Term count.} $H$ is a sum of exactly $\binom{N}{4}$ Pauli
strings: $495$ at $N=12$, $1820$ at $N=16$, $10\,626$ at $N=24$,
$35\,960$ at $N=32$ --- growing as $N^4/24$ while the Hilbert space grows as
$2^{N/2}$. For a fixed target precision $\epsilon$ the shot budget scales like
(sum of $|{\rm coefficients}|$)$^2/\epsilon^2$, so this is the obstruction that
bites first in practice.

Grouping helps, and the measured numbers say \emph{how much}: a greedy
(Welsh--Powell) colouring of the non-commutation graph packs the $1820$ strings
at $N=16$ into $167$ general-commuting groups but only $549$ qubit-wise
commuting ones. General commutation is therefore worth roughly a further
factor of three over the naive qubit-wise grouping, at the price of a Clifford
basis change per group.

\paragraph{(b) The terms are global.} The Jordan--Wigner strings make the mean
Pauli weight $5.1$ out of $8$ qubits at $N=16$, i.e.\ $O(n)$ rather than
$O(1)$. There is no locality to exploit: no light-cone argument, no local
cost-function trick, and every term must be measured on essentially the whole
register.

\paragraph{(c) The ground state is volume-law entangled.} This is the
obstruction that matters, and the one that cannot be optimised away.

\begin{center}\small
\input{tables/model3_ed}
\end{center}

\begin{figure}[h]\centering
\includegraphics[width=\textwidth]{figures/model3_obstructions.pdf}
\caption{(a) Pauli-string count versus number of simultaneously measurable
groups. (b) Half-cut entanglement entropy: the SYK ground state sits at
$76$--$93\%$ of the maximal value $\tfrac{n}{2}\ln 2$ and grows linearly in
$n$, whereas the XXX ground state grows logarithmically. (c) Variance of
$\partial E/\partial\theta_0$ at random initialisation, falling by a factor
$\sim\!2.6$ between $n=4$ and $n=8$: consistent with exponential decay, but
over too narrow a range to fit a slope one should believe (the dashed line is a
guide to the eye).}
\end{figure}

The measured half-cut entropy of the SYK ground state is $76$--$93\%$ of the
maximum $\tfrac{n}{2}\ln 2$ and grows \emph{linearly} in $n$; the XXX ground
state, by contrast, grows logarithmically. A circuit of depth $d$ with local
gates can only generate entanglement $O(d)$ across a cut, so reproducing a
near-maximally entangled state requires depth growing with $n$ --- not a
constant-depth, hardware-efficient ansatz. This is an expressivity statement,
independent of the optimiser, and it is why "add more layers" is not the
answer.

Two side observations from the same runs. The disorder fluctuation
$\sigma(E_0)/N$ falls from $2.4\times10^{-2}$ at $N=8$ to
$1.3\times10^{-3}$ at $N=20$ --- the model self-averages, so a handful of
realisations suffices and the error bars in the table are small on purpose, not
by luck. And the ground-state degeneracy alternates: in the range studied it is
$1$ for $N \equiv 0 \pmod 8$ ($N=8,16$) and $2$ otherwise, which is the known
dependence of the SYK spectrum on $N \bmod 8$ through the particle-hole
symmetry. It matters practically, because a VQE run over a degenerate ground
space converges to an arbitrary member of it, and comparing states rather than
energies then becomes meaningless.

\paragraph{(d) Flat gradients.}

\begin{center}\small
\input{tables/model3_bp}
\end{center}

The variance of a single energy derivative at random initialisation falls by a
factor $\sim\!2.6$ between $n=4$ and $n=8$. That is consistent with the
exponential decay expected of a barren plateau, but five points over four
qubits is too narrow a range to fit an exponent that anyone should believe, and
I would rather say so than quote a slope. The honest statement is that the
gradients are already small and getting smaller, and that this compounds with
the absence of any structure in the all-to-all random couplings for the
optimiser to follow.

\paragraph{(e) The ground state is complex, and the standard ansatz is not.}

This one is not in the usual list, and it turned out to matter more than (d).
The measured relative error of the textbook RY\,+\,CZ hardware-efficient
ansatz is \emph{identical to six digits} across the whole depth sweep $P=2,4,8$ --- $41.7\%$ at
$N=8$, $54.2\%$ at $N=10$ --- which is not what a barren plateau looks like.
It is what a representability failure looks like:

\begin{center}\small
\input{tables/model3_naive}
\end{center}

RY and CZ generate only real amplitudes. Measuring how complex the SYK ground
state actually is --- the residual norm of $\mathrm{Im}\,e^{-i\varphi}\psi_0$
after minimising over the global phase --- gives $0.49$ at $N=8$ and $0.66$ at
$N=10$, i.e.\ the state is about as far from real as a state can be. The
real-amplitude ansatz therefore saturates at the best real approximation, and
no depth can help. Before blaming the optimiser it is worth checking whether
the ansatz can represent the target at all; here it cannot, and that accounts
for essentially all of the $40$--$50\%$: simply adding one $R_Z$ per qubit per
layer makes the same brickwork \emph{exact} at $N=8$ and $N=10$.

That last observation also bounds what obstructions (c) and (d) can be blamed
for at these sizes. With $n=4$ or $5$ qubits the Hilbert space has $16$ or $32$
dimensions and a $40$-parameter circuit can cover it, so the binding constraint
in the range we can diagonalise is \emph{reality}, not entanglement. The
volume-law argument is an asymptotic statement about what a constant-depth
ansatz can do as $n$ grows; it is the reason not to expect the complex
brickwork to keep winning, not an explanation of the failure observed at
$N=8$. Keeping those two apart is the whole point of measuring them
separately.

\subsection{What I did about it: ADAPT-VQE}
\label{sec:adapt}

Rather than assert that an adaptive ansatz would do better, it is implemented.
The operator pool is the Hamiltonian's own quartic terms
$P_m = \chi_i\chi_j\chi_k\chi_l$, each a single Pauli string with $P_m^2=1$, so
\begin{equation}
  e^{-i\theta P_m} = \cos\theta\,\mathbb{1} - i\sin\theta\,P_m
\end{equation}
costs two sparse matrix--vector products. The operator added at each step is
the one with the largest
\begin{equation}
  \left.\frac{\partial E}{\partial \theta}\right|_{\theta=0}
  = \bra{\psi} i\,[P_m, H] \ket{\psi}
  = -2\,\mathrm{Im}\ip{P_m\psi}{H\psi},
\end{equation}
and it enters at $\theta=0$, so the growth is a warm start by construction ---
the same mechanism that gained an order of magnitude in Model 1.

\begin{center}\small
\input{tables/model3b}
\end{center}

The result is a negative one, and it is the most useful thing in this section.

\emph{ADAPT does not beat the complex brickwork at these sizes.} With a fixed
budget of $6n$ operators it reaches $12.1\%$ at $N=12$ against the brickwork's
$8.5\%$, and $27.1\%$ against $19.4\%$ at $N=14$. It is enormously better than
the \emph{real} brickwork, but that comparison flatters it, because the real
brickwork is failing for a trivial reason. Against an ansatz that can at least
represent a complex state, an adaptive construction with a Hamiltonian-derived
pool buys nothing here. I would rather report that than tune the budget until
the ordering flips.

\emph{What the sweep does show cleanly is where the expressivity wall is.} The
complex brickwork is \emph{exact} at $n=4$ and $n=5$, then degrades to $8.5\%$
at $n=6$ and $19.4\%$ at $n=7$, at fixed depth $P=4$. That is the onset
predicted by (c): as soon as the Hilbert space is larger than the circuit can
cover, a constant-depth ansatz starts losing ground, and it loses it fast. The
real brickwork, by contrast, sits at $41$--$54\%$ at \emph{every} size and
every depth, which is the signature of a constraint that has nothing to do with
system size at all. Two different failure modes, cleanly separated by having
measured both.

\emph{One instance-specific curiosity, reported because it cost me an hour.} At
$N=8$, ADAPT started from a real product reference stalls after a single
operator, at a point where every pool gradient vanishes to machine precision
($\sim10^{-16}$) --- the construction is trapped in the real subspace and the
pool, which generates the trap, cannot select its way out. A complex reference
(one $R_Z$ per qubit) unblocks it immediately. At $N=10,12,14$ the same real
reference does \emph{not} stall, so this is a property of that instance rather
than a general obstruction. It is worth knowing mostly as a diagnostic: a run
that dies after one operator with zero gradients is a symmetry, not a barren
plateau, and the fix is the reference state, not the optimiser.

\subsection{What I would do with more time}
\label{sec:more}

Ranked by what the measurements above support:

\begin{enumerate}\itemsep3pt
\item \textbf{Fix the ansatz before anything else.} Use a complex-capable
  circuit, and check the reference state. Those two changes take minutes and
  account for the entire difference between a $50\%$ error and an $8.5\%$ one
  at $N=12$; everything else in this list is a refinement on top.
\item \textbf{If keeping ADAPT, enlarge the pool beyond the Hamiltonian's own
  terms.} The $N=8$ stall shows that a pool sharing a symmetry with the
  reference cannot escape it, and the $N=12,14$ numbers show that a pool of
  quartic terms alone is not enough to beat a generic circuit.
\item \textbf{Measurement reduction.} General-commuting grouping
  ($167$ groups instead of $1820$ terms at $N=16$), then classical shadows,
  whose sample complexity depends on observable weight rather than term count
  and is therefore the natural tool for the global operators of (b).
\item \textbf{Warm start from a solvable neighbour} --- the free-fermion
  ($q=2$) Majorana model, whose ground state is Gaussian and preparable exactly
  in $O(n^2)$ gates. This is the only idea here that attacks (c) rather than
  working around it, because it builds the entanglement gradually instead of
  hoping to find it. \emph{It is implemented in Section~\ref{sec:warm}, and it
  is the strategy that works}: the same ADAPT pool and budget that gave
  $3$--$4\%$ from a product state reach the exact ground state at $N=10$ when
  started from the Gaussian state.
\item \textbf{Report the error curve, not the energy.} With $2^{N/2}=1024$ at
  $N=20$, exact diagonalisation is available far beyond where the VQE is
  useful, so the honest deliverable at these sizes is the gap between them.
\end{enumerate}

I would not claim that any of this makes SYK a good VQE target at accessible
$N$. It is a good \emph{stress test} precisely because it violates every
assumption --- locality, sparsity, low entanglement, reality --- that makes
Models 1 and 2 work.

\section{Software validation}\label{sec:soft}

The fast in-house state-vector engine (block-structured gate
\eqref{eq:gate}, adjoint gradients) is what makes $L=16$ affordable; Qiskit is
used to prove that it implements the intended circuit. At random parameters,
the two agree on the energy to $\sim5\times10^{-15}$ and on the state to
twelve digits, after accounting for the endianness reversal. The transpiled
circuit metrics below are what one would actually pay on hardware.

\begin{center}\small
\input{tables/qiskit_check}
\end{center}

\noindent Other checks performed: analytic gradient against central finite
differences ($\le 3\times10^{-9}$); $S^z=0$ sector against the full Hilbert
space for $L\le12$; the geometric form of $H_{\rm HS}$ against the
$\sin^2$ form, operator-by-operator; Hermiticity of every Hamiltonian;
$\{\chi_a,\chi_b\}=2\delta_{ab}$ by explicit matrices.

\section{Sampling cost, measured}
\label{sec:shots}

Every energy elsewhere in this report is an exact expectation value on the
state vector. That is deliberate --- the questions ask for the \emph{variational}
error, and shot noise would confound the two --- but it leaves the measurement
cost unquantified, which for a VQE is half the problem. So it is measured.

The estimator packs the Pauli strings into qubit-wise-commuting groups, rotates
each qubit into the basis its group requires, and draws bitstrings from
$|\psi|^2$; every member of a group is then estimated from the \emph{same}
shots, with their true correlations rather than an independence assumption.
It is verified unbiased against the exact value (\textsf{D.12}). The variance
scales as $1/M_g$, so the shots needed for a target precision follow from one
measurement at a known $M_g$.

\begin{center}\small
\input{tables/model4_shots}
\end{center}

Three things are worth reading off this table.

The ordering is \emph{not} the ordering of term counts. The XXX chain has the
fewest Pauli strings ($24$) and needs the most shots, because what enters the
budget is $\|c\|_1 = 24$, not the number of terms; Haldane--Shastry has three
and a half times as many strings but a coefficient norm of $9.7$, spread over
rapidly decaying $1/\sin^2$ couplings, and costs thirty times less.

Grouping pays where the structure allows it. Both spin chains collapse to
\textbf{three} qubit-wise-commuting groups --- all the $XX$ terms, all the $YY$,
all the $ZZ$ --- independently of $L$. The SYK model, whose strings are
essentially random, needs $41$ groups at $N=8$ and $196$ at $N=12$, and that
gap is exactly the price of having no structure to exploit.

And the absolute numbers are sobering: $\sim\!10^8$ shots for a single energy
evaluation of SYK at $N=12$ to $10^{-3}$, before any optimisation iteration.
Multiplied by the hundreds of iterations an optimiser needs, this --- not the
circuit depth --- is what makes the model hard in practice at accessible sizes.

\section{A Gaussian warm start, and the best result for SYK}
\label{sec:warm}

Of the strategies listed in Section~\ref{sec:more}, the $q=2$ warm start is
the cheapest to implement and turns out to be the only one that actually
works. It is therefore implemented rather than recommended.

The free-fermion Majorana model
\begin{equation}
  H_{q=2} = \frac{i}{2}\sum_{i<j} K_{ij}\,\chi_i\chi_j ,
  \qquad K \text{ real antisymmetric},
\end{equation}
is quadratic, so its ground state is Gaussian and preparable exactly on
hardware by an $O(n^2)$ network of Givens rotations. The explicit factor of $i$
is required here and not for the quartic terms: reversing a product of $q$
distinct anticommuting Majoranas costs $(-1)^{\binom q2}$, which is $+1$ for
$q=4$ but $-1$ for $q=2$ (\textsf{D.4}). Two uses were tested: as the reference
state of ADAPT, and as the $s=0$ end of the interpolation
$H(s) = (1-s)H_{q=2} + s\,H_{\rm SYK}$ with the complex brickwork re-optimised
along $s = 0, 0.25, 0.5, 0.75, 1$.

\begin{center}\small
\input{tables/model4_warmstart}
\end{center}

\begin{figure}[h]\centering
\includegraphics[width=0.56\textwidth]{figures/model4_strategies.pdf}
\caption{Every strategy tried on the SYK model, on one axis. The real-amplitude
brickwork (navy) is flat at $40$--$55\%$ because it cannot represent a complex
state at all. Allowing complex amplitudes (red) is exact at $n=4,5$ and then
degrades as the Hilbert space outgrows the circuit. ADAPT from a product state
(teal) is no better. ADAPT from the Gaussian $q=2$ ground state (gold) is exact
at $N=10$ and $0.13\%$ at $N=12$, and the advantage erodes by $N=14$.}
\end{figure}

Three readings, and the third is the one that matters most.

\emph{Overlap is not the mechanism.} The Gaussian state is a poor
approximation to the SYK ground state --- overlap $0.43$ at $N=10$, and only
$0.07$ at $N=12$ and $0.09$ at $N=14$. It belongs to a different model. What it
supplies is a starting point with the \emph{right kind} of correlations: a
fermionic Gaussian state already carries extensive entanglement, so the circuit
has to build corrections rather than build everything from a product state.
The size at which the overlap is smallest, $N=12$, is the size at which the
warm start helps most.

\emph{Used as the ADAPT reference it is decisive, at $N=12$ spectacularly so.}
The same pool, the same operator budget and the same optimiser that gave
$12.1\%$ from a product state give $0.13\%$ from the Gaussian state --- a factor
of $65$ --- and reach the exact ground state at $N=10$. This is the honest
resolution of the negative result of Section~\ref{sec:adapt}: ADAPT with a
Hamiltonian-derived pool was not the problem; starting it from a product state
was.

\emph{And the advantage erodes by $N=14$}, where ADAPT from the Gaussian state
gives $18.7\%$ against the cold brickwork's $19.4\%$ --- better, but no longer
by a distance. That is worth stating plainly rather than stopping at $N=12$.
The likely cause is the fixed operator budget of $6n$, which grows linearly
while the entanglement it has to represent grows linearly \emph{in the exponent};
the warm start buys a constant factor, not a change of scaling. Testing that
would mean growing the budget with $n$ and watching where the curves cross,
which is the obvious next run and is not in here.

The practical conclusion of the whole SYK section, which took three attempts to
find, is an ordering: \emph{fix the reference state before touching the ansatz,
and fix the ansatz before touching the optimiser.} Taken in the other order,
all one learns is that the optimiser has a hard time --- which is true, and
useless.

\section{Limitations}

Stated plainly, because they are the parts I would want to be asked about.
The $1/L^2$ extrapolation of Model 1 ignores the logarithmic corrections of the
critical chain, which is why it misses Bethe by $3\times10^{-3}$ rather than
by the fit error. The $L=16$ HVA runs use fewer restarts than the smaller
sizes, so the reported errors there are upper bounds, and the $L=16$ warm-start
ladder is the single most expensive computation here. The SYK exact
diagonalisation averages over $5$--$20$ disorder realisations, enough given the
measured self-averaging but not enough for tail statistics. The ADAPT runs use a pool truncated to the
largest-coefficient terms and a fixed operator budget of $6n$, so the reported
errors are again upper bounds; with a larger budget the numbers can only
improve. The shot-noise study of Section~\ref{sec:shots} uses a single
reference state per model (the exact ground state), so the variances quoted are
representative rather than worst case over the optimisation trajectory. The
$q=2$ warm start is tested on one disorder realisation per size and one choice
of $K_{ij}$ scaling; the $N=10$ result is exact and the trend is clear, but the
sample is small. And the optimiser throughout is L-BFGS-B or Powell on a
noiseless landscape --- a real device would need something noise-aware
(SPSA, or a surrogate model), which would change the iteration counts but none
of the conclusions about expressivity.

\vspace{0.4em}
\hrule
\vspace{0.4em}
\noindent\small
\textbf{Code.} \texttt{spinvqe/} contains the library
(\texttt{core.py} Hamiltonians and exact diagonalisation, \texttt{ansatz.py}
circuits and adjoint gradients, \texttt{vqe.py} the optimiser,
\texttt{syk.py} the Majorana mapping and Pauli bookkeeping,
\texttt{qiskit\_interface.py} the Qiskit export and cross-check).
\texttt{run\_model1.py}, \texttt{run\_model1\_warmstart.py},
\texttt{run\_model2.py}, \texttt{run\_model3.py} reproduce every number;
\texttt{make\_tables.py} and \texttt{make\_figures.py} regenerate the tables
and figures in this document from the JSON outputs in \texttt{results/}.

\appendix
\clearpage
\section{Derivations}
\label{app:math}

Everything used in the body is derived here, and every identity carries a tag
(\textsf{A.1}, \textsf{B.5}, \dots) that is also the name of a numerical check
in \texttt{verify\_derivations.py}. A derivation that is only written down can
be wrong in a way no reader notices; all $47$ of them are therefore evaluated
as well, and Table~\ref{tab:derivs} lists the residual of each.

\subsection{Conventions and the two-site algebra}

\paragraph{\textsf{A.1} Pauli versus spin.}
With $\bm\sigma = 2\bm S$,
\begin{equation}
  \ssig_i\cdot\ssig_j = 4\,\bm S_i\cdot\bm S_j .
  \tag{A.1}
\end{equation}
Every textbook result quoted below is stated for $\bm S\cdot\bm S$ and carries
this factor of $4$ when translated to the Hamiltonians as written in the
assignment. It is the single easiest way to be wrong by $4$ in this problem.

\paragraph{\textsf{A.2}--\textsf{A.3} Spectrum of the exchange operator.}
Let $P_{ij}$ be the operator that swaps sites $i$ and $j$. A one-line
computation on the four basis states gives
\begin{equation}
  P_{ij} = \tfrac12\big(\mathbb{1} + \ssig_i\cdot\ssig_j\big)
  \quad\Longleftrightarrow\quad
  \ssig_i\cdot\ssig_j = 2P_{ij} - \mathbb{1}.
  \tag{A.3}
\end{equation}
Since $P^2=\mathbb{1}$, $P$ has eigenvalues $\pm1$: $+1$ on the symmetric
(triplet) subspace and $-1$ on the antisymmetric (singlet) one. Hence
\begin{equation}
  \operatorname{spec}\big(\ssig\cdot\ssig\big) = \{\,+1 \text{ (triplet, 3-fold)},\;
  -3 \text{ (singlet)}\,\}.
  \tag{A.2}
\end{equation}
This single fact drives everything in Model 1: the gate, its cost, its
symmetry, and its differentiation rule.

\paragraph{\textsf{A.6} The gate in closed form.}
On $\{\ket{00},\ket{11}\}$, $\ssig\cdot\ssig$ acts as the identity; on
$\{\ket{01},\ket{10}\}$ its matrix is $\begin{psmallmatrix}-1&2\\2&-1\end{psmallmatrix}
= -\mathbb{1} + 2X$. Exponentiating each block,
\begin{equation}
  e^{-i\theta\,\ssig\cdot\ssig} =
  \begin{cases}
    e^{-i\theta} & \text{on } \ket{00},\ket{11},\\[2pt]
    e^{+i\theta}\big(\cos 2\theta\,\mathbb{1} - i\sin 2\theta\,X\big)
      & \text{on } \ket{01},\ket{10}.
  \end{cases}
  \tag{A.6}
\end{equation}
Applying this by slicing a reshaped state vector, rather than contracting a
generic $4\times4$ tensor, is what makes $L=16$ with $P=10$ layers run in
minutes rather than hours.

\paragraph{\textsf{A.4}--\textsf{A.5} Circuit form.}
On a fixed pair, $(XX)(YY) = (XY)\otimes(XY) = (iZ)\otimes(iZ) = -ZZ$ and
$(YY)(XX) = (-iZ)\otimes(-iZ) = -ZZ$ as well, so $XX$ and $YY$ commute; the
other two pairs are identical by symmetry. Commuting exponents factorise, so
\begin{equation}
  e^{-i\theta(XX+YY+ZZ)} = R_{XX}(2\theta)R_{YY}(2\theta)R_{ZZ}(2\theta),
  \qquad R_{PP}(\varphi) \equiv e^{-i\varphi PP/2}.
  \tag{A.5}
\end{equation}

\subsection{Model 1: ansatz, symmetry and gradients}

\paragraph{\textsf{B.0} The Bethe benchmark.}
For $H_S = J\sum \bm S_i\cdot\bm S_{i+1}$ the Bethe ansatz gives
$E/L \to J(\tfrac14 - \ln 2)$. With (A.1),
\begin{equation}
  \frac{E_0}{L} \longrightarrow 4J\left(\tfrac14-\ln 2\right)
  = J(1-4\ln 2) = -1.7725887\,J .
  \tag{B.0}
\end{equation}

\paragraph{No Trotter error.}
For even $L$ with periodic boundaries the bonds split into
$B_{\rm odd}=\{(0,1),(2,3),\dots\}$ and
$B_{\rm even}=\{(1,2),\dots,(L{-}1,0)\}$, each a set of \emph{disjoint} pairs.
Disjoint two-site operators commute, so
$e^{-i\theta H_{\rm odd}} = \prod_{(a,b)\in B_{\rm odd}} e^{-i\theta\,\ssig_a\cdot\ssig_b}$
\emph{exactly}. The circuit implements the ansatz with no splitting error, so
the only error in the tables is variational.

\paragraph{\textsf{B.1}--\textsf{B.3} Why the ansatz cannot leave the physical
sector.}
Each generator $\ssig_a\cdot\ssig_b$ is an $SU(2)$ scalar: writing
$\bm S_{\rm tot} = \sum_i \bm S_i$, one has
$[\bm S_{\rm tot}, \ssig_a\cdot\ssig_b]=0$ because the dot product of two
vector operators is invariant. Hence every gate --- and therefore the whole
circuit --- commutes with $\bm S_{\rm tot}^2$ and $S^z_{\rm tot}$. The
reference state is a product of two-site singlets, for which
$\bm S^2_{\rm tot}\ket{\psi_0}=0$ and $S^z_{\rm tot}\ket{\psi_0}=0$, so
\begin{equation}
  \bm S_{\rm tot}^2\ket{\psi(\bm\theta)} = 0,\qquad
  S^z_{\rm tot}\ket{\psi(\bm\theta)} = 0
  \quad\text{for all }\bm\theta .
  \tag{B.2--B.3}
\end{equation}
The antiferromagnetic ground state of the even-$L$ chain lives in exactly that
sector (Marshall's theorem), so no expressivity is wasted on states that are
orthogonal to the target by symmetry. Numerically, the amplitude of
$\ket{\psi(\bm\theta)}$ outside $S^z=0$ is $0$ to machine precision and
$\|\bm S^2_{\rm tot}\ket{\psi}\| \sim 10^{-15}$, at random $\bm\theta$.

\paragraph{\textsf{B.4} Adjoint gradients.}
Write $\ket{\psi} = U_M \cdots U_1\ket{\psi_0}$ with $U_k = e^{-i\theta_k H_k}$
and $\ket{\psi_k} = U_{k-1}\cdots U_1\ket{\psi_0}$. Since $H_k$ is a sum of
\emph{disjoint} terms, $\partial_{\theta_k}U_k = -iH_kU_k$, and
\begin{align}
  \frac{\partial E}{\partial\theta_k}
  &= 2\,\mathrm{Re}\,\bra{\psi}H\,U_M\cdots(\partial_{\theta_k}U_k)\cdots U_1\ket{\psi_0}
   \nonumber\\
  &= 2\,\mathrm{Re}\,\bra{\lambda_k}(-iH_k)\ket{\psi_{k+1}},
  \qquad
  \ket{\lambda_k} \equiv \Big(\prod_{j>k}U_j\Big)^{\!\dagger} H \ket{\psi}.
  \tag{B.4}
\end{align}
Both $\ket{\psi_{k+1}}$ and $\ket{\lambda_k}$ are obtained by sweeping once
forwards and once backwards, so the whole gradient costs about two forward
passes, independent of the number of parameters. Measured against central
differences: agreement $\le 3\times10^{-9}$, speed-up $\times 8$ at $L=16$.

\paragraph{\textsf{B.5}--\textsf{B.6} Parameter-shift rules, correctly.}
Let $U(\theta)=e^{-i\theta G}$. Expanding in the eigenbasis of $G$, the matrix
elements of $U^\dagger H U$ carry phases $e^{i(\lambda_j-\lambda_k)\theta}$, so
$E(\theta)$ is a trigonometric polynomial whose frequencies are the
\emph{distinct positive gaps} of $\operatorname{spec}(G)$.

For a \emph{single bond}, (A.2) gives gaps $\{0,4\}$: a single frequency
$\Omega=4$, hence $E(\theta)=a+b\cos4\theta+c\sin4\theta$ and the two-term rule
applies in rescaled form,
\begin{equation}
  \frac{dE}{d\theta}
  = \frac{\Omega}{2}\left[E\!\left(\theta+\frac{\pi}{2\Omega}\right)
    - E\!\left(\theta-\frac{\pi}{2\Omega}\right)\right]
  = 2\left[E\!\left(\theta+\tfrac{\pi}{8}\right)-E\!\left(\theta-\tfrac{\pi}{8}\right)\right].
  \tag{B.5}
\end{equation}
The natural guess is wrong in both directions: the \emph{naive} $\pi/2$-shift,
$\tfrac12$-prefactor version fails, while the rescaled one above is exact ---
verified to $3\times10^{-10}$.

For a whole sub-layer, $G=\sum_{r=1}^{R}(\ssig\cdot\ssig)_r$ over $R=L/2$
disjoint bonds. Its eigenvalues are $R-4k$, $k=0,\dots,R$, so the positive gaps
are $4,8,\dots,4R$: $R$ distinct frequencies, and an exact rule needs $2R$
evaluations. Numerically, at $L=8$ ($R=4$) the two-term rule is off by
$2.2$ in absolute terms --- not a small error, a wrong answer --- while the
$2R$-term rule is exact to $1.4\times10^{-9}$. With $2P$ parameters this would
cost $2P\cdot 2R = 2PL$ circuit evaluations per gradient against two passes
for (B.4), which is why the code uses (B.4).

\paragraph{\textsf{B.7} Why the $1/L^2$ extrapolation misses by $3\times10^{-3}$.}
The XXX chain is critical, described by a $c=1$ conformal field theory, for
which $E_0(L)/L = e_\infty - \pi v c/(6L^2) + \dots$. The subleading terms are
\emph{not} a clean series in $1/L^2$: the isotropic point carries a marginally
irrelevant operator that produces corrections of relative order
$1/\ln^3 L$. A two-parameter fit on five sizes therefore cannot land on
$e_\infty$ to better than about $10^{-3}$, and does not. Quoting the deviation
as a failure of the fit rather than as a physical feature would be the mistake.

\subsection{Model 2: Haldane--Shastry}

\paragraph{\textsf{C.1} The geometric identity.}
With $z_j=e^{2\pi i j/L}$,
\begin{align}
  |z_j-z_k|^2
  &= \big|e^{2\pi ij/L}-e^{2\pi ik/L}\big|^2
   = \big|e^{\pi i(j+k)/L}\big|^2\,
     \big|e^{\pi i(j-k)/L}-e^{-\pi i(j-k)/L}\big|^2 \nonumber\\
  &= 1\cdot\big|2i\sin\big(\pi(j-k)/L\big)\big|^2
   = 4\sin^2\!\Big(\frac{\pi(j-k)}{L}\Big).
  \tag{C.1}
\end{align}

\paragraph{\textsf{C.2} The Hamiltonian as written \emph{is} the textbook one.}
Substituting (C.1) and (A.1), the two factors of $4$ cancel:
\begin{equation}
  H_{\rm HS}
  = \frac{\pi^2 J}{L^2}\sum_{j<k}\frac{\ssig_j\cdot\ssig_k}{|z_j-z_k|^2}
  = \frac{\pi^2 J}{L^2}\sum_{j<k}\frac{4\,\bm S_j\cdot\bm S_k}
      {4\sin^2(\pi(j-k)/L)}
  = J\Big(\frac{\pi}{L}\Big)^{2}\sum_{j<k}
      \frac{\bm S_j\cdot\bm S_k}{\sin^2(\pi(j-k)/L)} .
  \tag{C.2}
\end{equation}
The complex-coordinate dressing of the problem statement is therefore
cosmetic --- but only once (C.1) is in hand.

\paragraph{\textsf{C.3} Closed-form energy.}
Given (C.2) the standard Haldane--Shastry result applies verbatim,
\begin{equation}
  E_0 = -\frac{\pi^2 J}{24}\Big(L+\frac{5}{L}\Big),
  \qquad \frac{E_0}{L}\xrightarrow[L\to\infty]{} -\frac{\pi^2}{24}
  = -0.4112335 .
  \tag{C.3}
\end{equation}
Verified against exact diagonalisation at $L=4,6,8,10,12,14$: agreement
$\le 9\times10^{-15}$ at every size. It is used as a benchmark only after that
check, not on authority.

\paragraph{\textsf{C.4}--\textsf{C.6} The exact ground state, and why the
momentum factor matters.}
In the $S^z=0$ sector, with down spins at sites $n_1<\dots<n_M$, $M=L/2$,
\begin{equation}
  \psi_{\rm J}(n_1,\dots,n_M)\;\propto\;
  \prod_{i=1}^{M} z_{n_i}\;\prod_{i<j}\big(z_{n_i}-z_{n_j}\big)^2 .
  \tag{C.4}
\end{equation}
The Jastrow factor $\prod_{i<j}(z_{n_i}-z_{n_j})^2$ is a squared Vandermonde
determinant --- the hallmark of the Gutzwiller projected Fermi sea, which is
what the Haldane--Shastry ground state is. The extra $\prod_i z_{n_i}$ is a
pure momentum factor: the lattice translation $T: z_n\mapsto z_{n+1}$
multiplies each coordinate by $e^{2\pi i/L}$, and it fixes the total momentum
to the value the ground state actually carries. Dropping it leaves a state that
is \emph{still an exact eigenvector} (residual $\sim10^{-15}$) but in a
different momentum sector: its overlap with the true ground state is zero to
machine precision, at $L=6,8,10$. That is the kind of detail which is invisible
in the algebra and immediate in the numerics, and the reason both variants were
implemented.

The deeper reason such a simple closed form exists is that the $1/r^2$ chain is
not merely integrable but carries a Yangian symmetry, so its spectrum is that
of free "spinons" with no scattering phase --- which is also why (C.3) is a
finite polynomial in $L$ rather than the solution of Bethe equations.

\paragraph{What this implies for a variational circuit.}
Model 2 has a zero-parameter exact ansatz, so running a VQE on it would be an
approximation to something already known. The two honest uses of (C.4) are as a
\emph{warm start} for a perturbed, non-integrable version of the model (say
$1/r^{2+\epsilon}$), where the circuit only has to represent corrections; and
as a benchmark for what happens when an ansatz built on nearest-neighbour bonds
is applied to an all-to-all Hamiltonian. Note also the change in measurement
cost: $H_{\rm HS}$ has $3\binom{L}{2}$ Pauli terms against $3L$ for the XXX
chain ($84$ against $24$ at $L=8$).

\subsection{Model 3: Jordan--Wigner and the SYK model}

\paragraph{\textsf{D.1}--\textsf{D.2} The mapping.}
With $n=N/2$ qubits and $k=0,\dots,n-1$,
\begin{equation}
  \chi_{2k} = \Big(\prod_{j<k}Z_j\Big)X_k,
  \qquad
  \chi_{2k+1} = \Big(\prod_{j<k}Z_j\Big)Y_k .
  \tag{D.0}
\end{equation}
Three cases establish $\{\chi_a,\chi_b\}=2\delta_{ab}$. \emph{(i)} $a=b$: the
$Z$-string squares to the identity and $X_k^2=Y_k^2=\mathbb{1}$, so
$\chi_a^2=\mathbb{1}$. \emph{(ii)} Same qubit, different operator:
$\{X_k,Y_k\}=0$ and the strings cancel. \emph{(iii)} Different qubits $k<l$:
the string of the operator with the larger index contains $Z_k$, and $Z_k$
anticommutes with both $X_k$ and $Y_k$, while every other factor commutes ---
so the pair anticommutes. Verified by explicit matrices for $N=4,\dots,12$.

\paragraph{\textsf{D.3}--\textsf{D.4} Each quartic term is one Hermitian Pauli
string.}
A product of Pauli strings is a Pauli string times a phase in
$\{\pm1,\pm i\}$, so $\chi_i\chi_j\chi_k\chi_l$ is a single string. Reversing a
product of $q$ distinct mutually anticommuting factors costs
$(-1)^{\binom{q}{2}}$; for $q=4$ this is $(-1)^6=+1$, so
\begin{equation}
  \big(\chi_i\chi_j\chi_k\chi_l\big)^\dagger
  = \chi_l\chi_k\chi_j\chi_i
  = \chi_i\chi_j\chi_k\chi_l ,
  \tag{D.4}
\end{equation}
and $H$ is Hermitian for real couplings, with no factor of $i$ needed. (For
$q=2$ the same count gives $(-1)^1=-1$, which is why the free-fermion model of
Section~\ref{sec:warm} carries an explicit $i$.)

\paragraph{\textsf{D.5}--\textsf{D.6} Counting.}
With $J_{ijkl}$ totally antisymmetric, the $4!$ orderings of each index set
contribute equally and cancel the $1/4!$ of the problem statement, so
\begin{equation}
  H_{\rm SYK} = \sum_{i<j<k<l} J_{ijkl}\,\chi_i\chi_j\chi_k\chi_l ,
  \qquad
  \overline{J_{ijkl}}=0,\quad \overline{J^2_{ijkl}}=\frac{6J^2}{N^3},
  \tag{D.5}
\end{equation}
which is exactly $\binom{N}{4}$ Pauli strings on $N/2$ qubits: $N^4/24$ terms
against a Hilbert space of $2^{N/2}$. The Jordan--Wigner $Z$-strings make the
typical string act on $O(n)$ qubits --- measured mean weight $5.1$ out of $8$
at $N=16$ --- so there is no locality to exploit.

\paragraph{\textsf{D.7} Commutation criterion.}
Two Pauli strings either commute or anticommute. On each site they
anticommute iff both factors are non-identity and different; the total sign is
the product over sites, so
\begin{equation}
  [A,B]=0 \iff
  \#\{\,k:\ A_k\neq I,\ B_k\neq I,\ A_k\neq B_k\,\}\ \text{is even}.
  \tag{D.7}
\end{equation}
This is the predicate coloured in the measurement-grouping analysis; at $N=16$
it packs $1820$ strings into $167$ groups, against $549$ under the stricter
qubit-wise criterion.

\paragraph{\textsf{D.8}--\textsf{D.9} The ADAPT gate and its selection
gradient.}
Each pool element is a Pauli string with $P^2=\mathbb{1}$, so
\begin{equation}
  e^{-i\theta P} = \cos\theta\,\mathbb{1} - i\sin\theta\,P ,
  \tag{D.8}
\end{equation}
two sparse matrix--vector products. Appending it at the end of the circuit,
$E(\theta)=\bra{\psi}e^{i\theta P}He^{-i\theta P}\ket{\psi}$ and
\begin{equation}
  \left.\frac{\partial E}{\partial\theta}\right|_{\theta=0}
  = \bra{\psi}i[P,H]\ket{\psi}
  = i\big(\ip{P\psi}{H\psi}-\overline{\ip{P\psi}{H\psi}}\big)
  = -2\,\mathrm{Im}\,\ip{P\psi}{H\psi},
  \tag{D.9}
\end{equation}
using $P^\dagger=P$ and $H^\dagger=H$. This is what ranks the pool; because
each new operator enters at $\theta=0$, the construction is a warm start by
design.

\paragraph{\textsf{D.12} Shot cost.}
For $H=\sum_s c_s P_s$ measured in groups $g$ of mutually commuting strings,
each single shot of group $g$ returns an unbiased estimate of every member
simultaneously, and
\begin{equation}
  \operatorname{Var}\big[\hat E\big]
  = \sum_g \frac{1}{M_g}
    \operatorname{Var}_g\Big[\sum_{s\in g}c_s P_s\Big]
  \;\le\; \frac{\big(\sum_s|c_s|\big)^2}{M},
  \tag{D.12}
\end{equation}
so the worst-case budget for precision $\epsilon$ is
$(\sum_s |c_s|)^2/\epsilon^2$. The implementation samples bitstrings from
$|\psi|^2$ after rotating each qubit into the required basis, so the
correlations within a group are the true ones rather than an assumption; the
estimator is verified unbiased against the exact expectation value.

\paragraph{The reality obstruction, stated precisely.}
A circuit built from $R_Y$ and $CZ$ acting on $\ket{0}^{\otimes n}$ has real
amplitudes in the computational basis, so it can only represent states with
$\mathrm{Im}\,e^{-i\varphi}\psi = 0$ for some global phase $\varphi$. Defining
\begin{equation}
  \kappa(\psi) \;=\; \min_{\varphi}\;
  \frac{\big\|\mathrm{Im}\,e^{-i\varphi}\psi\big\|}{\|\psi\|}
  \;\in\;[0,1],
\end{equation}
a real-amplitude ansatz is fundamentally limited whenever $\kappa>0$. For the
SYK ground state we measure $\kappa = 0.49$ at $N=8$ and $0.66$ at $N=10$,
which is why the $R_Y+CZ$ brickwork returns the \emph{same} energy at every
depth, and why one $R_Z$ per qubit per layer makes it exact at those sizes.

\paragraph{The entanglement obstruction, stated precisely.}
A circuit of depth $d$ built from two-qubit gates can only act across a given
cut $d$ times, so the entanglement entropy it can generate across that cut is
$O(d)$. The SYK ground state has half-cut entropy $76$--$93\%$ of the maximum
$\tfrac n2\ln 2$, growing linearly in $n$; the XXX ground state grows
logarithmically. Hence a constant-depth ansatz cannot scale here --- a
statement about expressivity, independent of the optimiser. At the sizes where
we can diagonalise ($n\le 7$) this is not yet the binding constraint, and
saying otherwise would be overclaiming: it becomes visible as the complex
brickwork degrades from exact at $n=5$ to $8.5\%$ at $n=6$ and $19.4\%$ at
$n=7$.

\subsection{The verification table}

\begin{center}\scriptsize
\input{tables/derivations}
\captionof{table}{Every identity of this appendix, evaluated numerically by
\texttt{verify\_derivations.py}. $47/47$ pass. Entry \textsf{B.6a} is inverted
by construction: it records that the two-term shift rule \emph{fails} for a
four-bond sub-layer, with an observed error of $2.2$.}
\label{tab:derivs}
\end{center}


\end{document}
